\section{Memory}

\begin{remark}
    Memory is the \textbf{limiting factor} in computation. It limits computation in time, uses by far most of the space and consumes most of the energy.
\end{remark}

\begin{remark}
    The robustness decreases with increased memory capacity. (Ex. Row-Hammering-Attack)
\end{remark}

\begin{remark}
    What the programmer sees is \textbf{virtual memory addresses} that are mapped to \textbf{physical memory} behind the scenes. This simplifies programming.
\end{remark}

\begin{definition}
    \textbf{Ideal memory} has zero latency, cost, energy and infinite capacity and bandwidth
\end{definition}

\begin{definition}
    \textbf{Memory types} are defined by the trade-offs of cost and performance.
    \begin{itemize}
        \item FlipFlops and Latches: (tens of transistors per bit)
        \item Static Random Access Memory (SRAM): (6 transistors per bit)
        \item Dynamic Random Access Memory (DRAM): (cheaper but slower, needs refreshing)
        \item Storage: (slow but cheap, 1 transistor per bit, non-volatile)
    \end{itemize}
\end{definition}

\subsection{organization of Memory}

\begin{definition}
    Memory is an array of cells with adressing and output logic.
\end{definition}

\begin{remark}
    To increase performance memory is divided into multiple banks that allow concurrent access to independent adresses.
\end{remark}

\begin{definition}
    A typical \textbf{DRAM hierarchy} has
    \begin{itemize}
        \item Multiple-channels that leave/enter the memory controller 
        \item Dual in-line memory modules (DIMM) which are mounted on the motherboard
        \item Ranks most likely front/back of card which are independently adressable bank groups with that share command and data busses.
        \item Chips that have multiple banks and share data bus out of the chip
        \item Bank can be made of multiple mats and has a global row/collumn organization
    \end{itemize}
\end{definition}

\begin{definition}
    A bank hase a \textbf{row-buffer} where a whole row is loaded and from where the column bit is selected.
\end{definition}

\begin{remark}
    If a row is already in the row-buffer it is a row hit, else a row-miss.
\end{remark}

\begin{remark}
    To load a row into the row-buffer we use a row of sense amplifiers that sense the bit in the capacitors, amplify and then restore them.
\end{remark}

\begin{remark}
    On a high level,SRAM is organized similarly but uses more transistors and is faster.
\end{remark}

\subsection{Memory Technologies}

\begin{definition}
    \textbf{Static Random Access Memory (SRAM)} is a low latency memory built from two cross-coupled inverters that uses 6 transistors and does not need any refresh cycles.
\end{definition}

\begin{definition}
    \textbf{Dynamic Random Access Memory (DRAM)} uses only 1 transistor and a capacitor to store a bit. It is cheaper but slower and needs to be refreshed periodically.
\end{definition}

\begin{definition}
    \textbf{Phase change memory (PCM)} is a new non-volatile memory technology that uses chalcogenide glass in its Amorphous and Crystalline forms to store bits. The form of the glass can be changed by applying heat using current to the material. 
\end{definition}

\begin{remark}
    PCM is slow, energy hungry and has limited endurance (the material degrades over time). However, it is non-volatile and very dense.
\end{remark}

\subsection{Memory Hierarchy}

\begin{remark}
    The reason to having multiple levels of memory are the trade-offs between cost, capacity, latency and energy consumption. We can not have fast and large memory at the same time hence we have a small fast memory (caches) and a large slow memory (RAM). 
\end{remark}

\begin{remark}
    By intelligently selecting and moving chunks of data form the main memory to the caches we get nearly the performance of a large and fast memory.
\end{remark}

\begin{definition}
   \textbf{Data Locality} is the property of data that allows us to predict where it will be accessed next.
   \begin{itemize}
       \item \textbf{Temporal locality:} If a data item is accessed, it is likely to be accessed again soon.
       \item \textbf{Spatial locality:} If a data item is accessed, it is likely that nearby data items will be accessed soon.
   \end{itemize}
\end{definition}

\begin{definition}
    \textbf{Manually managed memory} is when the programmer has to manually tell the system when to put what into which cache. These are often refered to as shared memory or scratchpads.
\end{definition}
